\subsection{What the Julia delegates preserve, and what the model must identify}
\label{rel22:survival-delegates}

The survival proposal becomes more precise when the two named delegates are
read alongside their caller. The selected SurvivalBase revision declares
version 0.1.0 in June 2026; SurvivalDistributions declares 0.1.1 in July 2026.
Both carry direct MIT notices. These are pre-cutoff source snapshots, not an
observed dependency resolution or a tested installation of SurvivalModels.
Their narrower roles are valuable: SurvivalBase constructs formula terms,
and SurvivalDistributions supplies distribution families and hazard operations.
Neither supplies the treatment assignment, censoring independence or
transport assumptions needed for an intervention recommendation.
\citep{rel22survbase,rel22survdist,rel22survcaller}

SurvivalBase's \texttt{Surv} term produces two columns: observed time and an
event indicator. Its README explicitly limits the current binding to
right-censored observations. The selected term has no entry-time field or
pair of interval-censoring endpoints. The parent general-hazard likelihood
then uses the event contribution minus cumulative hazard from origin. This
composition is coherent for its stated observation type, but passing two
convenient columns cannot encode a delayed-entry experiment. Under independent
entry and censoring, conditional on the modeled covariates, an observation
entering at time $l$ and exiting at time $t$ instead contributes
\begin{equation}
 \ell=\delta\log h(t)-H(t)+H(l),\qquad t>l,
 \label{rel22:entry-likelihood}
\end{equation}
where $h$ is the hazard, $H$ its cumulative integral and $\delta$ indicates an
observed event. The last term conditions on survival to entry. It is not
recovered by relabeling entry as an interval-censoring endpoint or by retaining
only the exit time. \citep{rel22survbase,rel22survcaller,rel22survwitness}

An exact small construction shows the consequence for risk costing. Two
accounts enter at month four; one defaults at month five and the other is
right censored at month six. In a stipulated exponential model with constant
rate $\lambda$, the conditional exposure is three account-months and the
maximum-likelihood rate is $1/3$. Treating both accounts as observed from
origin gives eleven account-months and a rate of $1/11$. For a currently
surviving account, an assumed USD~100 loss on default gives next-month
expected losses of USD~28.35 and USD~8.69, respectively, using
$100(1-e^{-\lambda})$. These are calculations under the constructed
observation law, not observed bank losses or fitted Julia outputs. They show
why a right-censored worker must reject unsupported entry information or
delegate to a likelihood that retains it; an engine's convenient schema
cannot decide the scientific clock. \citep{rel22survwitness}

Even correctly encoded observations need an identifiable parameterization.
The parent general-hazard implementation composes the baseline as
\begin{equation}
 h(t\mid x,z)=h_0(te^{z\alpha})e^{x\beta},\qquad
 H(t\mid x,z)=H_0(te^{z\alpha})e^{x\beta-z\alpha}.
 \label{rel22:gh-composition}
\end{equation}
The second expression follows by changing variables in the hazard integral.
The caller's time and hazard multipliers implement this distinction; simply
multiplying cumulative hazard by the hazard-level factor would omit the
time-change correction. Yet retaining both coefficient blocks in an output
does not prove that the data distinguish them. Take a Weibull baseline with
shape $k>0$, scale $s>0$ and one covariate shared by the two blocks, so $z=x$.
Then
\begin{equation}
 H(t\mid x)=(t/s)^k
       \exp\!\left(x\{(k-1)\alpha+\beta\}\right).
 \label{rel22:weibull-ridge}
\end{equation}
Only the displayed linear combination enters the conditional law. Replacing
the coefficient pair by
$(\alpha+c,\beta-(k-1)c)$ leaves every hazard, survival probability and
right-censored likelihood contribution unchanged. For shape two, the pairs
$(0,\log 3)$ and $(\log 2,\log(3/2))$ are an exact example. One labels the
covariate as a pure hazard-level change; the other assigns part to the time
scale, although they predict the same event distribution.
\citep{rel22survcaller,rel22survwitness}

The parent removes coefficient blocks that are inactive by method definition
before inverting an observed Hessian. That avoids one source of singularity,
but its general-hazard branch retains both blocks and does not remove this
baseline-dependent ridge. With an exponential baseline, an accelerated-hazard
time coefficient disappears altogether because the baseline hazard is
constant. A finite optimizer result, inverse-Hessian request or table headed
\emph{identified parameters} cannot establish separate identification. The
defensible alternatives are to restrict the scientific model, report the
identified combination and its uncertainty, or justify a different baseline
and covariate design with diagnostics for weak identification. No distribution
fit alone turns either coefficient into a causal effect of payment holidays.
These statements concern exact parameter redundancy; no optimizer failure
frequency is claimed. \citep{rel22survcaller,rel22survwitness}

Parameter domains also have to survive the call boundary. The selected
LogLogistic family represents a logistic distribution for log event time.
Writing its location as $\mu\in\mathbb R$ and its positive scale as $\sigma$,
its median is $e^\mu$. The parent's baseline-fitting path exponentiates every
baseline parameter, however, including this location. For a baseline-only fit
without regression covariates, a median of half a month requires
$\mu=\log(1/2)<0$ and is outside that transformed fitting domain. Expressing
the same law in days gives a median of fifteen and location $\log 15>0$.
Changing units should not change whether a physical distribution belongs to
the admitted model. A distribution-specific transformation must distinguish
unrestricted locations from positive scales and shapes. Regression may
absorb some scale changes in other designs, so this example is deliberately
limited to the baseline-only fit. It neither asserts that every baseline is
affected nor predicts the optimizer's returned approximation.
\citep{rel22survdist,rel22survcaller,rel22survwitness}

\subsection{Tail behavior and the meaning of a survival output}
\label{rel22:survival-tail}

The piecewise-constant hazard delegate exposes a second consequential
boundary. It declares that the last supplied hazard continues after the final
finite time interval. Its inspected quantile loop nevertheless returns that
interval's endpoint when the requested cumulative hazard exceeds the
accumulated finite-grid hazard. The cumulative-hazard tail refers to
\texttt{L} although the argument is \texttt{D}; the expectation method takes
\texttt{L} but accesses \texttt{D}. Those names are not bound by the inspected
module. The abstract hazard fallback likewise calls numerical helpers that
the selected module does not import or define. These are specific source
obstacles to admitting these paths, not evidence that every distribution in
the package fails. Fixing a variable name alone would not establish the
intended tail law or qualify the numerical implementation.
\citep{rel22survdist}

For a one-month grid interval with hazard $1/5$ per month continuing forever,
the stipulated law is simply exponential. Its ninetieth percentile is
\begin{equation}
 q_{0.9}=5\log 10\simeq11.513\text{ months},
 \label{rel22:tail-quantile}
\end{equation}
whereas the finite-loop trace returns one month. Survival at that endpoint
is still about 0.819, so the difference cannot be dismissed as a negligible
tail. At a stipulated continuous receipt rate of USD~100 per month until
termination, with no discounting, competing event or cost, expected total
receipts are USD~500. Receipts restricted to the first month are only
USD~90.63. Both are legitimate targets when labeled; silently imposing the
grid endpoint would substitute one for the other. A bank requiring a finite
valuation horizon should specify that horizon and integrate to it. A model
claiming an unbounded mean must establish a proper tail and finite expectation.
The example reconstructs the source branch algebraically and does not run the
downloaded Julia implementation. \citep{rel22survdist,rel22survwitness}

Output names are another source of silent substitution. In the parent,
\texttt{predict\_expected} returns cumulative hazard; the survival method
exponentiates its negative. Cumulative hazard is neither a mean event time nor
an event probability. If it equals two at the chosen horizon, the corresponding
event probability is $1-e^{-2}\simeq0.865$. With an assumed USD~100 loss on
that event, expected loss is USD~86.47, not USD~200. The no-time prediction
also evaluates each training subject at its own observed time, whereas a
portfolio decision generally requires a common future horizon. The adapter
must carry target, units, population and horizon with each array rather than
infer them from a method name. \citep{rel22survcaller,rel22survwitness}

These findings support a conditional comparison, with distinct admission
decisions for right-censored parametric fitting, a piecewise hazard
representation and a general-hazard coefficient model. Preserve the useful
formula and distribution interfaces while testing entry/censoring meaning,
parameter domains, common-horizon prediction, likelihood rank, tail inversion
and time-unit changes against independent laws. Include optimizer convergence
and uncertainty failures in the result: the inspected fit stores the minimizer
and objective without retaining a convergence flag in its model fields.
Native qualification and transitive dependency review remain to be done.
The direct MIT notices identify the selected components' declared terms;
they do not establish scientific calibration or clear every dependency.
For the banking decision, choose the least elaborate model whose assumptions
and identified outputs answer the question, and retain uncertainty about the
causal interpretation separately from uncertainty within a fitted survival law.
\citep{rel22survbase,rel22survdist,rel22survcaller}

